\documentclass[12pt,a4paper]{article}
\usepackage[utf8]{inputenc}
\usepackage[T1]{fontenc}
\usepackage{amsmath}
\usepackage{amssymb}
\usepackage{amsfonts}
\usepackage{mathtools}
\usepackage{physics}
\usepackage{siunitx}
\usepackage{graphicx}
\usepackage{xcolor}
\usepackage{hyperref}
\usepackage{geometry}
\usepackage{microtype}
\usepackage{lmodern}
\geometry{margin=2.5cm}
\hypersetup{colorlinks=true,linkcolor=blue,urlcolor=blue}

\title{JFM   LaTeX Template for Journal of Fluid Mechanics.pdf}
\author{Jack Ma}
\date{07 mai 2026}

\begin{document}
\maketitle

\section*{JFM   LaTeX Template for Journal of Fluid Mechanics.pdf}
Banner appropriate to article type will appear here in typeset article

\title{
Weak Solutions of Conservation Laws: An Application to Inviscid Burger's Equation and Shallow-Water Equations
}

\author{
Makimona Kiakisolako \({ }^{1}\) \\ \({ }^{1}\) Department of Mathematics, African Institute for Mathematical Sciences, Kigali, Rwanda \\ Corresponding author: bienaime.kiakisolako@aims.ac.rw
}
(Received xx; revised xx; accepted xx)

Modeling natural phenomena is an important problem; however, it can be difficult due to factors that can create discontinuities making it difficult to find the expected solutions to the models. Hence, we need to look for other types of solution with less restrictive conditions on the differentiability of the solutions to the models.

As uniqueness under these less restrictive conditions is not guaranteed, conditions from gas dynamics are used to choose the more physically realistic solution as the true one. We illustrate the notions we discuss in an example of a real fluid model called the shallow water equations.

Key words: Geophysical and Geological Flows

\section*{1. First-order heading}

The layout design for the Journal of Fluid Mechanics journal has been implemented as a LaTeX style file.The FLM style file is based on the ARTICLE style as discussed in the LaTeX manual. Commands which differ from the standard LaTeX interface, or which are provided in addition to the standard interface, are explained in this guide. This guide is not a substitute for the LaTeX manual itself.

\subsection*{1.1. Introduction to LaTeX}

The LaTeX document preparation system is a special version of the TeX typesetting program. LaTeX adds to TeX a collection of commands which simplify typesetting by allowing the author to concentrate on the logical structure of the document rather than its visual layout.
LaTeX provides a consistent and comprehensive document preparation interface. There are simple-to-use commands for generating a table of contents, lists of figures and/or tables, and indexes. LaTeX can automatically number list entries, equations, figures, tables, and

Author 1 et al

\begin{figure}
\includegraphics[alt={},max width=\textwidth]{https://cdn.mathpix.com/cropped/69b520a6-acfb-49ed-a568-5ce682ad3f27-02.jpg?height=504&width=1139&top_left_y=182&top_left_x=311}
\captionsetup{labelformat=empty}
\caption{Figure 1. This is a sample figure caption extended to multiple rows. This is a sample figure caption extended to multiple rows. This is a sample figure caption extended to multiple rows.}
\end{figure}

\begin{table}
\begin{tabular}{lccc}
\(a / d\) & \(M=4\) & \(M=8\) & Callan et al. \\
0.1 & 1.56905 & 1.56 & 1.56904 \\
0.3 & 1.50484 & 1.504 & 1.50484 \\
0.55 & 1.39128 & 1.391 & 1.39131 \\
0.7 & 1.32281 & 10.322 & 1.32288 \\
0.913 & 1.34479 & 100.351 & 1.35185
\end{tabular}
\captionsetup{labelformat=empty}
\caption{Table 1. This is a sample table caption}
\end{table}
footnotes, as well as parts, chapters, sections and subsections. Using this numbering system, bibliographic citations, page references and cross references to any other numbered entity (e.g. chapter, section, equation, figure, list entry) are quite straightforward.

\subsection*{1.2. The FLM document class}

The use of document class allows a simple change of style (or style option) to transform the appearance of your document. The CUP FLM class file preserves the standard LaTeX interface such that any document which can be produced using the standard LaTeX ARTICLE style can also be produced with the FLM style. However, the fonts (sizes) and measure of text is slightly different from that for ARTICLE, therefore line breaks will change and it is possible that equations may need re-setting.

\section*{2. Figures and Tables}

\subsection*{2.1. Figures}

Each figure should be accompanied by a single caption, to appear beneath, and must be cited in the text. Figures should appear in the order in which they are first mentioned in the text. For example see figures 1 and 2.

\subsection*{2.2. Tables}

Tables, however small, must be numbered sequentially in the order in which they are mentioned in the text. Words table 1, table 2 should be lower case throughout. See table 1 for an example.

\begin{figure}
\includegraphics[alt={},max width=\textwidth]{https://cdn.mathpix.com/cropped/69b520a6-acfb-49ed-a568-5ce682ad3f27-03.jpg?height=788&width=1232&top_left_y=191&top_left_x=216}
\captionsetup{labelformat=empty}
\caption{Figure 2. This is a sample figure caption with (a) and (b) designations to define parts.}
\end{figure}

\section*{3. Notation and style}

Generally any queries concerning notation and journal style can be answered by viewing recent pages in the Journal. However, the following guide provides the key points to note. It is expected that Journal style and mathematical notation will be followed, and authors should take care to define all variables or entities upon first use. Also note that footnotes are not normally accepted. Abbreviations must be defined at first use, glossaries or lists/tables of abbreviations are not permitted.

\subsection*{3.1. Mathematical notation}

\subsection*{3.1.1. Setting variables, functions, vectors, matrices etc}
- Italic font should be used for denoting variables, with multiple-letter symbols avoided except in the case of dimensionless numbers such as \(R e, P r\) and \(P e\) (Reynolds, Prandtl, and Péclet numbers respectively, which are defined as \Rey, \Pran and \(\backslash\) Pen in the template).
- Upright Roman font (or upright Greek where appropriate) should be used for:
1. (vI) label, e.g. T. t (transpose)
2. Fixed operators: sin, log, d, \(\Delta, \exp\) etc.
3. Constants: \(\mathrm{i}(\sqrt{-1}), \pi\) (defined as \upi),e etc.
4. Special Functions: Ai, Bi (Airy functions, defined as \Ai and \Bi), Re (real part, defined as \Real), Im (imaginary part, defined as \Imag), etc.
5. Physical units: \(\mathrm{cm}, \mathrm{s}\), etc.
6. Abbreviations: c.c. (complex conjugate), h.o.t. (higher-order terms), DNS, etc.
- Bold italic font (or bold sloping Greek) should be used for vectors (with the centred dot for a scalar product also in bold): \(\boldsymbol{i} \cdot \boldsymbol{j}\)
- Bold sloping sans serif font, defined by the \mathsfbi macro, should be used for tensors and matrices: D
- Calligraphic font (for example \(\mathcal{G}, \mathcal{R}\) ) can be used as an alternative to italic when the same letter denotes a different quantity use \mathcal in \(\mathrm{IAT}_{\mathrm{E}} \mathrm{X}\) )

\subsection*{3.1.2. Other symbols}

Large numbers that are not scientific powers should not include commas, but should use a non-breaking space, and use the form 1600 or 16000 or 160000 . Use \(O\) to denote 'of the order of', not the LATEX \(O\).

The product symbol ( × ) should only be used to denote multiplication where an equation is broken over more than one line, to denote a cross product, or between numbers . The symbol should not be used, except to denote a scalar product of vectors specifically.

\subsection*{3.1.3. Example Equations}

This section contains sample equations in the JFM style. Please refer to the ETEX source file for examples of how to display such equations in your manuscript.
\[
\left.\begin{array}{c}
\left(\nabla^{2}+k^{2}\right) G_{s}=\left(\nabla^{2}+k^{2}\right) G_{a}=0 \\
\boldsymbol{\nabla} \cdot \boldsymbol{v}=0, \quad \nabla^{2} P=\boldsymbol{\nabla} \cdot(\boldsymbol{v} \times \boldsymbol{w}) \\
G_{s}, G_{a} \sim 1 /(2 \pi) \ln r \quad \text { as } \quad r \equiv|P-Q| \rightarrow 0 \\
\frac{\partial G_{s}}{\partial y}=0 \quad \text { on } \quad y=0 \\
G_{a}=0 \quad \text { on } \quad y=0
\end{array}\right\}, y \cos k(x-\xi) t \mathrm{~d} t, \quad 0<y, \quad \eta<d, y \begin{gathered}
-\frac{1}{2 \pi} \int_{0}^{\infty} \gamma^{-1}[\exp (-k \gamma|y-\eta|)+\exp (-k \gamma(2 d-y-\eta))] \cos (t)=\left\{\begin{array}{cl}
-\mathrm{i}\left(1-t^{2}\right)^{1 / 2}, & t \leq 1 \\
\left(t^{2}-1\right)^{1 / 2}, & t>1
\end{array}\right. \\
-\frac{1}{2 \pi} \int_{0}^{\infty} B(t) \frac{\cosh k \gamma(d-y)}{\gamma \sinh k \gamma d} \cos k(x-\xi) t \mathrm{~d} t \\
G=-\frac{1}{4} \mathrm{i}\left(H_{0}(k r)+H_{0}\left(k r_{1}\right)\right)-\frac{1}{\pi} \int_{0}^{\infty} \frac{\mathrm{e}^{-k \gamma d}}{\gamma \sinh k \gamma d} \cosh k \gamma(d-y) \cosh k \gamma(d-\eta)
\end{gathered}
\]

Note that when equations are included in definitions, it may be suitable to render them in line, rather than in the equation environment: \(\boldsymbol{n}_{q}=\left(-y^{\prime}(\theta), x^{\prime}(\theta)\right) / w(\theta)\). Now \(G_{a}=\frac{1}{4} Y_{0}(k r)+\widetilde{G_{a}}\) where \(r=\left\{[x(\theta)-x(\psi)]^{2}+[y(\theta)-y(\psi)]^{2}\right\}^{1 / 2}\) and \(\widetilde{G_{a}}\) is regular as \(k r \rightarrow 0\). However, any fractions displayed like this, other than \(\frac{1}{2}\) or \(\frac{1}{4}\), must be written on the line, and not stacked (ie 1/3).
\[
\begin{aligned}
\frac{\partial}{\partial n_{q}}\left(\frac{1}{4} Y_{0}(k r)\right) & \sim \frac{1}{4 \pi w^{3}(\theta)}\left[x^{\prime \prime}(\theta) y^{\prime}(\theta)-y^{\prime \prime}(\theta) x^{\prime}(\theta)\right] \\
& =\frac{1}{4 \pi w^{3}(\theta)}\left[\rho^{\prime}(\theta) \rho^{\prime \prime}(\theta)-\rho^{2}(\theta)-2 \rho^{\prime 2}(\theta)\right] \quad \text { as } \quad k r \rightarrow 0 \\
& \frac{1}{2} \phi_{i}=\frac{\pi}{M} \sum_{j=1}^{M} \phi_{j} K_{i j}^{a} w_{j}, \quad i=1, \ldots, M,
\end{aligned}
\]
where
\[
\begin{gathered}
K_{i j}^{a}= \begin{cases}\partial G_{a}\left(\theta_{i}, \theta_{j}\right) / \partial n_{q}, & i \neq j \\
\partial \widetilde{G_{a}}\left(\theta_{i}, \theta_{i}\right) / \partial n_{q}+\left[\rho_{i}^{\prime} \rho_{i}^{\prime \prime}-\rho_{i}^{2}-2 \rho_{i}^{\prime 2}\right] / 4 \pi w_{i}^{3}, & i=j \\
\rho_{l}=\lim _{\zeta \rightarrow Z_{l}^{-}(x)} \rho(x, \zeta), \quad \rho_{u}=\lim _{\zeta \rightarrow Z_{u}^{+}(x)} \rho(x, \zeta)\end{cases} \\
\left(\rho(x, \zeta), \phi_{\zeta \zeta}(x, \zeta)\right)=\left(\rho_{0}, N_{0}\right) \quad \text { for } \quad Z_{l}(x)<\zeta<Z_{u}(x) . \\
\tau_{i j}=\left(\overline{\bar{u}_{i} \bar{u}_{j}}-\bar{u}_{i} \bar{u}_{j}\right)+\left(\overline{\bar{u}_{i} u_{j}^{S G S}+u_{i}^{S G S} \bar{u}_{j}}\right)+\overline{u_{i}^{S G S} u_{j}^{S G S}}, \\
\tau_{j}^{\theta}=\left(\overline{\bar{u}_{j} \bar{\theta}}-\bar{u}_{j} \bar{\theta}\right)+\left(\overline{\bar{u}_{j} \theta^{S G S}+u_{j}^{S G S} \bar{\theta}}\right)+\overline{u_{j}^{S G S} \theta^{S G S}} . \\
\boldsymbol{Q}_{C}=\left[\begin{array}{ccccc}
-\omega^{-2} V_{w}^{\prime} & -\left(\alpha^{t} \omega\right)^{-1} & 0 & 0 & 0 \\
\frac{\beta}{\alpha \omega^{2}} V_{w}^{\prime} & 0 & 0 & 0 & \mathrm{i} \omega^{-1} \\
\mathrm{i} \omega^{-1} & 0 & 0 & 0 & 0 \\
\mathrm{i} R_{\delta}^{-1}\left(\alpha^{t}+\omega^{-1} V_{w}^{\prime \prime}\right) & 0 & -\left(\mathrm{i} \alpha^{t} R_{\delta}\right)^{-1} & 0 & 0 \\
\frac{\mathrm{i} \beta}{\alpha \omega} R_{\delta}^{-1} V_{w}^{\prime \prime} & 0 & 0 & 0 & 0 \\
\left(\mathrm{i} \alpha^{t}\right)^{-1} V_{w}^{\prime} & \left(3 R_{\delta}^{-1}+c^{t}\left(\mathrm{i} \alpha^{t}\right)^{-1}\right) & 0 & -\left(\alpha^{t}\right)^{-2} R_{\delta}^{-1} & 0
\end{array}\right] .
\end{gathered}
\]
where \(\hat{\boldsymbol{\eta}}^{t}=\boldsymbol{b} \exp \left(\mathrm{i} \gamma x_{3}^{t}\right)\).
\[
\begin{gathered}
\operatorname{Det}\left[\rho \omega^{2} \delta_{p s}-C_{p q r s}^{t} k_{q}^{t} k_{r}^{t}\right]=0 \\
\left\langle k_{1}^{t}, k_{2}^{t}, k_{3}^{t}\right\rangle=\left\langle\alpha^{t}, 0, \gamma\right\rangle \\
f(\theta, \psi)=(g(\psi) \cos \theta, g(\psi) \sin \theta, f(\psi)) \\
f\left(\psi_{1}\right)=\frac{3 b}{\pi\left[2\left(a+b \cos \psi_{1}\right)\right]^{3 / 2}} \int_{0}^{2 \pi} \frac{\left(\sin \psi_{1}-\sin \psi\right)(a+b \cos \psi)^{1 / 2}}{\left[1-\cos \left(\psi_{1}-\psi\right)\right](2+\alpha)^{1 / 2}} \mathrm{~d} x \\
g\left(\psi_{1}\right)=\frac{3}{\pi\left[2\left(a+b \cos \psi_{1}\right)\right]^{3 / 2}} \int_{0}^{2 \pi}\left(\frac{a+b \cos \psi}{2+\alpha}\right)^{1 / 2}\left\{f(\psi)\left[\left(\cos \psi_{1}-b \beta_{1}\right) S+\beta_{1} P\right]\right. \\
\times \frac{\sin \psi_{1}-\sin \psi}{1-\cos \left(\psi_{1}-\psi\right)}+g(\psi)\left[\left(2+\alpha-\frac{\left(\sin \psi_{1}-\sin \psi\right)^{2}}{1-\cos \left(\psi-\psi_{1}\right)}-b^{2} \gamma\right) S\right. \\
\left.\left.+\left(b^{2} \cos \psi_{1} \gamma-\frac{a}{b} \alpha\right) F\left(\frac{1}{2} \pi, \delta\right)-(2+\alpha) \cos \psi_{1} E\left(\frac{1}{2} \pi, \delta\right)\right]\right\} \mathrm{d} \psi \\
\alpha=\alpha\left(\psi, \psi_{1}\right)=\frac{b^{2}\left[1-\cos \left(\psi-\psi_{1}\right)\right]}{(a+b \cos \psi)\left(a+b \cos \psi_{1}\right)}, \quad \beta-\beta\left(\psi, \psi_{1}\right)=\frac{1-\cos \left(\psi-\psi_{1}\right)}{a+b \cos \psi} \\
H(0)=\frac{\epsilon \bar{C}_{v}}{\tilde{v}_{T}^{1 / 2}(1-\beta)}, \quad H^{\prime}(0)=-1+\epsilon^{2 / 3} \bar{C}_{u}+\epsilon \hat{C}_{u}^{\prime} \\
H^{\prime \prime}(0)=\frac{\epsilon u_{*}^{2}}{\tilde{v}_{T}^{1 / 2} u_{P}^{2}}, \quad H^{\prime}(\infty)=0
\end{gathered}
\]

Lemma 1. Let \(f(z)\) be a trial Batchelor (1971, pp. 231-232) function defined on \([0,1]\). Let \(\Lambda_{1}\) denote the ground-state eigenvalue for \(-\mathrm{d}^{2} g / \mathrm{d} z^{2}=\Lambda g\), where \(g\) must satisfy \(\pm \mathrm{d} g / \mathrm{d} z+\alpha g=0\) at \(z=0,1\) for some non-negative constant \(\alpha\). Then for any \(f\) that is not identically zero we have
\[
\frac{\alpha\left(f^{2}(0)+f^{2}(1)\right)+\int_{0}^{1}\left(\frac{\mathrm{~d} f}{\mathrm{~d} z}\right)^{2} \mathrm{~d} z}{\int_{0}^{1} f^{2} \mathrm{~d} z} \geq \Lambda_{1} \geq\left(\frac{-\alpha+\left(\alpha^{2}+8 \pi^{2} \alpha\right)^{1 / 2}}{4 \pi}\right)^{2} .
\]

Corollary 1. Any non-zero trial function \(f\) which satisfies the boundary condition \(f(0)=f(1)=0\) always satisfies
\[
\int_{0}^{1}\left(\frac{\mathrm{~d} f}{\mathrm{~d} z}\right)^{2} \mathrm{~d} z .
\]

\section*{4. Additional facilities}

In addition to all the standard LaTeX design elements, the FLM style includes the following feature:
- Extended commands for specifying a short version of the title and author(s) for the running headlines.
Once you have used this additional facility in your document, do not process it with a standard LaTeX style file.

\subsection*{4.1. Titles authors' names and affiliation}

In the FLM style, the title of the article and the author's name (or authors' names) are used both at the beginning of the article for the main title and throughout the article as running headlines at the top of every page. The Journal title is used on odd-numbered pages (rectos) and the author's name appears on even-numbered pages (versos). Although the main heading can run to several lines of text, the running head line must be a single line.

Moreover, the main heading can also incorporate new line commands (e.g. \(\backslash \backslash\) ) but these are not acceptable in a running headline. To enable you to specify an alternative short title and author's name, the standard \righttitle and \lefttitle commands have been used to print the running headline. \corresau\{\} command should be used to provide the corresponding author details as shown below.
\lefttitle\{A.N. Jones, H.-C. Smith and J.Q. Long\} \righttitle\{Journal of Fluid Mechanics\} \title\{JFM \{\LaTeX\} submission template A framework for assessing the Reynolds analogy\} \author\{Alan N. Jones \(\backslash a f f\{1\}\), H.-C. Smith\aff\{1\} \and J.Q. Long\aff\{2\}\} \affiliation\{ \(a f f\{1\}\) STM Journals, Cambridge University Press, The Printing House, Shaftesbury Road, Cambridge CB2 8BS, UK \aff\{2\}DAMTP, Centre for Mathematical Sciences, Wilberforce Road, Cambridge CB3 OWA, UK\} \corresau\{Alan N. Jones, \email\{Jones@univ.edu\}\}

\subsection*{4.2. Abstract}

The FLM style provides for an abstract which is produced by the following commands
\begin\{abstract\} ... \end\{abstract\} }

\subsection*{4.3. Keywords}

The FLM style provides for an keywords which is produced by the following commands
\begin\{keywords\}... \begin\{keywords\} }
4.4. Lists

The FLM style provides the three standard list environments.
- Bulleted lists, created using the itemize environment.
- Numbered lists, created using the enumerate environment.
- Labelled lists, created using the description environment.

\subsection*{4.5. Footnotes}

The FLM journal style uses superior numbers for footnote references. \({ }^{1}\)

\section*{5. Some guidelines for using standard facilities}

The following notes may help you achieve the best effects with the FLM style file.

\subsection*{5.1. Sections}

LaTeX provides five levels of section headings and they are all defined in the FLM style file:
- \section.
- \subsection.
- \subsubsection.
- \paragraph.
- \subparagraph.

Section numbers are given for sections, subsection and subsubsection headings.

\subsection*{5.2. Running headlines}

As described above, the title of the journal and the author's name (or authors' names) are used as running headlines at the top of every page. The title is used on odd-numbered pages (rectos) and the author's name appears on even-numbered pages (versos).

The \pagestyle and \thispagestyle commands should not be used. Similarly, the commands \markright and \markboth should not be necessary.

\subsection*{5.3. Illustrations (or figures)}

The FLM style will cope with most positioning of your illustrations and you should not normally use the optional positional qualifiers on the figure environment which would override these decisions. Figure captions should be below the figure itself, therefore the \caption command should appear after the figure or space left for an illustration.

Figure 3 shows an example on working with LaTeX code to load art files. \includegraphics commnad is to load art files scale option used in \includegraphics is to reduce the

\footnotetext{
\({ }^{1}\) This shows how a footnote is typeset.
}

\begin{figure}
\includegraphics[alt={},max width=\textwidth]{https://cdn.mathpix.com/cropped/69b520a6-acfb-49ed-a568-5ce682ad3f27-08.jpg?height=364&width=449&top_left_y=186&top_left_x=659}
\captionsetup{labelformat=empty}
\caption{Figure 3. An example figure with space for artwork.}
\end{figure}
art. EPS format will be compiled using LaTeX. Also, PNG, PDF and JPG format art files are loaded in the same command but the TeX file should be compiled using PDFLaTeX:
```
\begin{figure}
    \includegraphics[scale=.4]{sample.eps}
    \caption{An example figure with space for artwork.}
    \label{sample-figure}
\end{figure}
```


The vertical depth should correspond roughly to the artwork you will submit; it will be adjusted to fit the final artwork exactly.

\subsection*{5.4. Creating new theorem-like environments}

You can create your own environments in LaTeX, and although you may already be familiar with \newtheorem, you will not have seen the other two commands explained below.
\newtheorem is a standard command used for creating new theorem-like environments, such as theorems, corollaries, lemmas, conjectures and propositions, with the body of the text (automatically) in italic.

\section*{6. List of packages used in the template}

Below are the list of packages that are already used in template, so we don't need to copy these packages again in the TeX file.
- \usepackage\{etex\}
- \usepackage\{amsthm\}
- \usepackage\{amssymb\}
- \usepackage\{soul\}
- \usepackage\{calc\}
- \usepackage\{color\}
- \usepackage\{colortbl\}
- \usepackage[boxed]\{algorithm2e\}
- \usepackage\{epstopdf\}
- \usepackage\{booktabs\}
- \usepackage\{natbib\}
- \usepackage\{hyperref\}
- \usepackage\{breakurl\}
- \usepackage\{bookmark\}
- \usepackage\{graphicx\}
- \usepackage\{caption\}
- \usepackage\{newtxtext\}

0 X0-8
- \usepackage\{newtxmath\}

\section*{7. Mathematics}

The FLM class file will centre displayed mathematics, and will insert the correct space above and below if standard LaTeX commands are used; for example use \(\backslash[\ldots\).\(] ] and\) not \(\$ \$ \ldots . \$ \$\). Do not leave blank lines above and below displayed equations unless a new paragraph is really intended.
amsmath. sty is common package to handle various type math equations was used in template. The amsmath descriptions are available in the document can be find in the web link https://ctan.org/pkg/amsmath?lang=en

\subsection*{7.1. Numbering of equations}

The subequations and subeqnarray environments have been incorporated into the FLM class file (see Section 7.1.1 regarding the subequations environment). Using these two environments, you can number your equations (7.1a), (7.1b) etc. automatically. For example, you can typeset
\[
a_{1} \equiv\left(2 \Omega M^{2} / x\right)^{\frac{1}{4}} y^{\frac{1}{2}}
\]
and
\[
a_{2} \equiv(x / 2 \Omega)^{\text {TeXtstyle } \frac{1}{2}} k_{y} / M .
\]
by using the subequations environment as follows:
```
\begin{subequations}
\begin{equation}
    a_1 \equiv (2\Omega M^2/x)^{\textstyle\frac{1}{4}}
        y^{\textstyle\frac{1}{2}}\label{a1}
\end{equation}
and
\begin{equation}
    a_2 \equiv (x/2\Omega)^{\textstyle\frac{1}{2}}k_y/M.\label{a2}
\end{equation}
\end{subequations}
```


\subsection*{7.1.1. The subequations environment and the AMSTEX package}

The amstex (and the amsmath) packages also define a subequations environment. The environment in JFM-FLM_Au.cls is used by default, as the environments in the AMS packages don't produce the correct style of output.

Note that the subequations environment from the amstex package takes an argument - you should use an 'a' to give \alph style subequations. e.g.
\begin\{subequations\}\{a\} ... \end\{subequations\} }

\subsection*{7.2. Bibliography}

As with standard LaTeX, there are two ways of producing a bibliography; either by compiling a list of references by hand (using a thebibliography environment), or by using BibTeX with a suitable bibliographic database with the bibliography style provided with this FLMguide.tex like \bibliographystyle\{jfm\}. The "jfm.bst" will produce the bibliography which is similar to FLM style but not exactly. If any modification has to be made with "jfm.bst" can be adjusted during manuscript preparation but the updated bst file
should be given with source files. However, contributors are encouraged to format their list of references style outlined in section 7.2.2 below.

\subsection*{7.2.1. References in the text}

References in the text are given by author and date. Whichever method is used to produce the bibliography, the references in the text are done in the same way. Each bibliographical entry has a key, which is assigned by the author and used to refer to that entry in the text. There is one form of citation - \cite\{key\} - to produce the author and date. Thus, ? is produced by
\cite\{Arntzenius2012\}.
In FLM, for references natbib. sty is used. natbib. sty is common package to handle various reference and its cross citations. There different type of cross citation such as \citep, \citet, \citeyear etc. of the natbib descriptions are available in the document can be find in the web link https://ctan.org/pkg/natbib?lang=en

Sample of basic cross citations examples from natbib (?) and ?. Similarly other command can be utilized from referring the description from https://ctan.org/pkg/natbib? lang=en

If citations have to sort then use the class option "citesort".

\subsection*{7.2.2. List of references}

The following listing shows some references prepared in the style of the journal. \begin\{thebibliography\}\{\} } \bibitem[Batchelor (1971)]\{Batchelor59\} \{\sc Batchelor, G.K.\} 1971 \{Small-scale variation of convected quantities like temperature in turbulent fluid part1, general discussion and the case of small conductivity\}, \{\it J. Fluid Mech.\}, \{\bf 5\}, pp. 3-113-133.
\bibitem [Bouguet (2008)]\{Bouguet01\} \{\sc Bouguet, J.-Y\} 2008 Camera Calibration Toolbox for Matlab \{\url\{http://www.vision.caltech.edu/bouguetj/calib_doc/\}\}. \bibitem[Briukhanovetal et al (1967)] \{Briukhanovetal1967\} \{ Sc Briukhanov, A. V., Grigorian, S. S., Miagkov, S. M., Plam, M. Y., I. E. Shurova, I. E., Eglit, M. E. and Yakimov, Y. L.\} 1967 \{On some new approaches to the dynamics of snow avalanches\}, \{\it Physics of Snow and Ice, Proceedings of the International Conference on Low Temperature Science\} \{Vol 1\} pp. 1221--1241 \{Institute of Low Temperature Science, Hokkaido University, Sapporo, Hokkaido, Japan\}. \bibitem[Brownell (2004)]\{Brownell04\} \{\sc Brownell, C.J. and Su, L.K.\} 2004 \{Planar measurements of differential diffusion in turbulent jets\}, \{\it AIAA Paper\}, pp. 2004-2335. \bibitem[Brownell and Su (2007)] \{Brownell07\} \{\sc Brownell, C.J. and Su, L.K.\} 2007 \{Scale relations and spatial spectra in a differentially diffusing jet\}, \{\it AIAA Paper\}, pp 2007-1314.
\bibitem [Dennis (1985)] \{Dennis85\}\} \{\sc Dennis, S.C.R.\} 1985 \{Compact explicit finite difference
approximations to the Navier--Stokes equation\}, \{ In \it Ninth Intl Conf. on Numerical Methods in Fluid Dynamics\}, \{ed Soubbaramayer and J.P. Boujot\}, \{Vol 218\}, \{\it Lecture Notes in Physics\}, pp. 23-51. Springer.
\bibitem [Edwards et al. (2017)]\{EdwardsVirouletKokelaarGray2017\} \{ sc Edwards, A. N., Viroulet, S., Kokelaar, B. P. and Gray, J. M. N. T.\} 2017 Formation of levees, troughs and elevated channels by avalanches on erodible slopes \{\it J. Fluid Mech.\}, \{\bf 823\}, pp. 278-315.
\bibitem[Hwang et al (1970)] \{Hwang70\}
\{\sc Hwang, L.-S. and Tuck, E.O.\} 1970 On the oscillations of
harbours of arbitrary shape \{\it J. \({ }^{\sim}\) Fluid Mech.\}, \(\{\mathrm{b} £ 42\}\),
pp 447-464.
\bibitem[Josep and Saut (1990)] \{JosephSaut1990\}
\{\sc Joseph, Daniel D. and Saut, Jean Claude\} 1990 Short-wave
instabilities and ill-posed initial-value problems \{\it Theoretical
and Computational Fluid Dynamics\}, \{\bf 1\}, pp.191--227,
\{|url\{http://dx.doi.org/10.1007/BF00418002\}\}.
\bibitem[Worster (1992)] \{Worster92\}
\{\sc Worster, M.G.\} 1992 The dynamics of mushy layers \{\it Interactive
dynamics of convection and solidification\}, \{(ed. S.H. Davis and H.E.
Huppert and W. Muller and M.G. Worster)\}, pp. 113--138 \{Kluwer\}.
\bibitem[Koch(1983)] \{Koch83\}
\{\sc Koch, W.\} 1983 Resonant acoustic frequencies of flat plate
cascades \{\it J. ~́Sound Vib.\}, \{\bf 88\}, pp. 233-242.
\bibitem[Lee(1971)] \{Lee71\}
\{\sc Lee, J.-J.\} 1971 Wave-induced oscillations in harbours of
arbitrary geometry \{\it J. \({ }^{\sim}\) Fluid Mech.\}, \{ bf 45\}, pp. 375-394.
\bibitem[Linton and Evans (1992)] \{Linton92\}
\{ sc Linton, C.M. and Evans, D.V.\} 1992 The radiation and scattering
of surface waves by a vertical circular cylinder in a channel
\{\it Phil. T Trans. V R. Soc. \Lond.\}, \{\bf 338\}, pp. 325-357.
\bibitem [Martin(1980] \{Martin80\}
\{\sc Martin, P.A.\} 1980 On the null-field equations for the exterior
problems of acoustics \{\it Q. \({ }^{\sim}\) J. Mech. \Appl. Maths \(\},\{\backslash \mathrm{bf} \mathrm{33} \mathrm{\}}\),
pp. 385--396.
\bibitem [Rogallo(1981)] \{Rogallo81\}
    \{\sc Rogallo, R.S.\} 1981 Numerical experiments in homogeneous
    turbulence \{ \{\it Tech. Rep.\} 81835\} \{NASA Tech. \Mem\}.
\bibitem[Ursell(1950)] \{Ursell50\}
\{ sc Ursell, F.\} 1950 Surface waves on deep water in the presence of a
submerged cylinder i \{\it Proc.\Camb. \Phil. \Soc.\}, \{\bf 46\},
pp.141--152.
\bibitem[Wijngaarden (1968)]\{Wijngaarden68\}
\{ Sc van Wijngaarden, L.\} 1968 On the oscillations Near and at resonance
in open pipes \{\it J. \({ }^{\sim}\) Engng Maths\}, \(\{\backslash \mathrm{bf} 2\}, \mathrm{pp} .225--240\).
\bibitem[Miller (1991)]\{Miller91\}
\{\sc Miller, P.L.\} 1991 Mixing in high Schmidt number turbulent jets
\{school \{PhD thesis\}\} \{California Institute of Technology\}.
\end\{thebibliography\} }

This list typesets as shown at the end of this guide. Each entry takes the form
\bibitem[\protect\citename\{Author(s), \}Date]\{tag\}
Bibliography entry
where Author(s) should be the author names as they are cited in the text, Date is the date to be cited in the text, and tag is the tag that is to be used as an argument for the \cite\{\} command. Bibliography entry should be the material that is to appear in the bibliography, suitably formatted. This rather unwieldy scheme makes up for the lack of an author-date system in LaTeX.

\section*{8. Notes for Editors}

This appendix contains additional information which may be useful to those who are involved with the final production stages of an article. Authors, who are generally not typesetting the final pages in the journal's typeface (Monotype Times), do not need this information.

\subsection*{8.1. Editing reference citations}

There different type of cross citation such as \citep, \citet, \citeyear etc. of the natbib descriptions are available in the document can be find in the web link https: //ctan.org/pkg/natbib?lang=en.

Please use the exact natbib command to display reference citations like (?) "(Author et al., 1990)" use \citep\{key\} to get the desired output.

\section*{9. Citations and references}

All papers included in the References section must be cited in the article, and vice versa. Citations should be included as, for example "It has been shown (Rogallo 1981) that..." (using the \citep command, part of the natbib package) "recent work by Dennis (1985)..." (using \citet). The natbib package can be used to generate citation variations, as shown below.
\citet[pp. 2-4]\{Hwang70\}:
Hwang et al (1970, pp. 2-4)
\citep[p. 6]\{Worster92\}:
(Worster 1992, p. 6)
\citep[see][]\{Koch83, Lee71, Linton92\}:
(see Koch 1983; Lee 1971; Linton and Evans 1992)
\citep[see][p. 18]\{Martin80\}:
(see Martin 1980(@, p. 18)
\citep\{Brownell04,Brownell07,Ursell50,Wijngaarden68,Miller91\}:
(Brownell 2004; Brownell and Su 2007; Ursell 1950; Wijngaarden 1968; Miller 1991)
(Briukhanovetal et al 1967)
Bouguet (2008)
(Josep and Saut 1990)
The References section can either be built from individual \bibitem commands, or can be built using BibTex. The BibTex files used to generate the references in this document can be found in the JFM LATEX template files folder provided on the website here.

\section*{Journal of Fluid Mechanics}

Where there are up to ten authors, all authors' names should be given in the reference list. Where there are more than ten authors, only the first name should appear, followed by et al.

\section*{10. Miscellaneous section heads}

Philosophy of Science asks authors to include Acknowledgments, Declarations (of competing interests), and Funding Information in your typeset manuscript. Please add three sections, reflecting each of these categories, using "bmhead" coding as shown below.
\begin\{bmhead\}[Xxxxxxx.] }
For the custom heading such as acknowledgment, funding disclosure, conflict disclosure and any other like-wise sections must be mentioned in the optional braces as shown in this example. \end\{bmhead\}\} }
The output of the above coding is shown below:
Xxxxxxxx. For the custom heading such as acknowledgment, funding disclosure, conflict disclosure and any other like-wise sections must be mentioned in the optional braces as shown in this example.

\section*{Appendix A}

In order not to disrupt the narrative flow, purely technical material may be included in the appendices. This material should corroborate or add to the main result and be essential for the understanding of the paper. It should be a small proportion of the paper and must not be longer than the paper itself.

\section*{References}

Batchelor, G.K. 1971 Small-scale variation of convected quantities like temperature in turbulent fluid part1, general discussion and the case of small conductivity, J. Fluid Mech., 5, pp. 3-113-133.
Bouguet, J.-Y 2008 Camera Calibration Toolbox for Matlab http://www.vision.caltech.edu/ bouguetj/calib_doc/.
Briukhanov, A. V., Grigorian, S. S., Miagmov, S. M., Plam, M. Y., I. E. Shurova, I. E., Eglit, M. E. and Yakimov, Y. L. 1967 On some new approaches to the dynamics of snow avalanches, Physics of Snow and Ice, Proceedings of the International Conference on Low Temperature Science Vol 1 pp. 1221-1241 Institute of Low Temperature Science, Hokkaido University, Sapporo, Hokkaido, Japan.
Brownell, C.J. and Su, L.K. 2004 Planar measurements of differential diffusion in turbulent jets, AIAA Paper, pp. 2004-2335.
Brownell, C.J. and Su, L.K. 2007 Scale relations and spatial spectra in a differentially diffusing jet, AIAA Paper, pp 2007-1314.
Dennis, S.C.R. 1985 Compact explicit finite difference approximations to the Navier-Stokes equation, In Ninth Intl Conf. on Numerical Methods in Fluid Dynamics, ed Soubbaramayer and J.P. Boujot, Vol 218, Lecture Notes in Physics, pp. 23-51. Springer.
Edwards, A. N., Viroulet, S., Kokelaar, B. P. and Gray, J. M. N. T. 2017 Formation of levees, troughs and elevated channels by avalanches on erodible slopes J. Fluid Mech., 823, pp. 278-315.
Hwang, L.-S. and Tuck, E.O. 1970 On the oscillations of harbours of arbitrary shape J. Fluid Mech., 42, pp 447-464.
Joseph, Daniel D. and Saut, Jean Claude 1990 Short-wave instabilities and ill-posed initial-value problems Theoretical and Computational Fluid Dynamics, 1, pp.191-227, http://dx.doi.org/10. 1007/BF00418002.
Worster, M.G. 1992 The dynamics of mushy layers Interactive dynamics of convection and solidification, (ed. S.H. Davis and H.E. Huppert and W. Muller and M.G. Worster), pp. 113-138 Kluwer.

Koch, W. 1983 Resonant acoustic frequencies of flat plate cascades J. Sound Vib., 88, pp. 233-242.
Lee, J.-J. 1971 Wave-induced oscillations in harbours of arbitrary geometry J. Fluid Mech., 45, pp. 375-394.
Linton, C.M. and Evans, D.V. 1992 The radiation and scattering of surface waves by a vertical circular cylinder in a channel Phil. Trans. R. Soc. Lond., 338, pp. 325-357.
Martin, P.A. 1980 On the null-field equations for the exterior problems of acoustics Q. J. Mech. Appl. Maths,33, pp. 385-396.
Rogallo, R.S. 1981 Numerical experiments in homogeneous turbulence Tech. Rep. 81835 NASA Tech. Mem.
Ursell, F. 1950 Surface waves on deep water in the presence of a submerged cylinder i Proc. Camb. Phil. Soc., 46, pp.141-152.
van Wijngaarden, L. 1968 On the oscillations Near and at resonance in open pipes J. Engng Maths,2, pp. 225-240.
Miller, P.L. 1991 Mixing in high Schmidt number turbulent jets school PhD thesis California Institute of Technology.


\end{document}
